\chapter[chap:phase6]{Phase 6: Agenten-Schnittstellen}

\section{Ziel}

Agenten arbeiten über definierte Werkzeuge. Jeder Aufruf ist autorisiert und protokolliert.

\section{Umfang}

\startitemize
\item API-Server auf dem Kern, mit denselben Operationen wie die CLI.
\item MCP-Server als Adapter, der die API als Werkzeuge bereitstellt.
\item Protokollierung jedes Aufrufs mit Akteur, Modellversion und Session.
\item Begrenzung der Aufrufe pro Akteur.
\stopitemize

\section{Offene Fragen}

\startitemize
\item Auswahl der Werkzeuge: Wie grob oder fein sind die Operationen?
\item Authentifizierung und Ablauf von Zugriffstoken für Agenten.
\item Verhalten bei Konflikten zwischen parallelen Agenten.
\item Commits und Merges auf dem Hub für Agenten ohne Klon.
Phase 5 baut sie nicht, legt aber die Richtung fest, die Phase 6 beim Bauen bestätigen muss
(ADR-015): Der Hub rechnet und liefert, der Akteur setzt den Commit zusammen und signiert.
Auf Anfrage würde der Hub dann den Tree des sauberen Merges und eine Vorschau des Urteils des
Merge-Gates berechnen; der Agent baut den Merge-Commit selbst, lädt ihn hoch und signiert ein
gewöhnliches Ref-Update (\in{Kapitel}[chap:phase5]).
Der Hub erzeugt also keinen Commit und signiert nie für einen Akteur.
Die Grenzen: Ein Akteur ohne Klon hat kein Operation Log und damit keinen Anker, die Aufzeichnung
des Hubs ist die einzige Spur seiner Handlungen, und er vertraut der Vorschau.
Offen sind die Anfragen selbst, auch die für das Lesen ohne Klon, das an eine Referenz gebunden
bleibt.
\stopitemize
